\documentclass[11pt,a4paper]{article}
\usepackage[utf8]{inputenc}
\usepackage[T1]{fontenc}
\usepackage{lmodern}
\usepackage[margin=24mm]{geometry}
\usepackage{amsmath,amssymb,mathtools}
\usepackage{booktabs}
\usepackage{enumitem}
\usepackage{xcolor}
\usepackage[hidelinks]{hyperref}

\newcommand{\dd}{\mathop{}\!\mathrm{d}}
\newcommand{\ii}{\mathrm{i}}
\newcommand{\id}{\mathrm{id}}
\newcommand{\R}{\mathbb{R}}
\newcommand{\Z}{\mathbb{Z}}
\newcommand{\N}{\mathbb{N}}
\newcommand{\Ord}{\operatorname{ord}}
\newcommand{\Ker}{\operatorname{Ker}}
\newcommand{\Aut}{\operatorname{Aut}}
\newcommand{\slashed}[1]{\not\!#1}
\setlength{\parindent}{0pt}
\setlength{\parskip}{6pt}
\allowdisplaybreaks

\title{Particle Physics --- Interactions and Fields}
\author{IMAPP lecture notes}
\date{Spring 2026}
\begin{document}\maketitle

\section{Classical potential picture}
For a static source, the potential satisfies a Klein-Gordon-type equation,
\[
\left(\nabla^2-\frac{1}{R_0^2}\right)V(\mathbf R)=\delta^{(3)}(\mathbf R).
\]
For a massless mediator ($R_0\to\infty$), this gives the Coulomb form $V(R)\propto1/R$. For a massive mediator, the potential is Yukawa-like:
\[
V(R)\propto\frac{e^{-R/R_0}}{4\pi R},\qquad R_0\simeq\frac{1}{m_X}.
\]
The finite range of the strong nuclear force, roughly $R_0\sim1\,\mathrm{fm}$, corresponds to a mediator mass of order $200\,\mathrm{MeV}$.

\section{Exchange of a virtual particle}
At higher energies, interactions are described by quantum field theory. A scattering amplitude contains the couplings at the two vertices and the propagator of the exchanged particle:
\[
\mathcal M\sim g_a\,\frac{1}{q^2-m_X^2}\,g_b,
\]
where $q$ is the transferred four-momentum. The exchanged particle is virtual and need not satisfy $q^2=m_X^2$; it is off shell.

\section{Feynman channels}
For $2\to2$ scattering, the momentum transfer can be organized into channels:
\begin{itemize}
  \item $t$-channel exchange: $q=p_1-p_3$, typically spacelike ($q^2<0$).
  \item $s$-channel exchange: the incoming particles form a time-like intermediate state ($q^2=s>0$).
  \item $u$-channel exchange: the crossed version of the process.
\end{itemize}

\section{Interaction strengths}
The electromagnetic interaction is controlled by the fine-structure constant $\alpha\simeq1/137$. Strong interactions are characterized by colour charge and a scale-dependent coupling. Weak interactions involve flavour and/or neutrino changes and are suppressed at low energy by massive-boson propagation, with a typical factor proportional to $1/M_W^2$.

\section*{Summary}
The potential picture becomes a field-theoretic exchange picture: range is set by mediator mass, amplitudes contain propagators, and scattering is classified by $s$, $t$, and $u$ channels.
\end{document}
